\documentclass[../../../include/open-logic-section]{subfiles}

\begin{document}

\olfileid{sth}{spine}{recursion}
\olsection{مبرهنة (مبرهنات) العودية المتناهية التجاوز}

يلزمنا أولًا أن نثبت أن~\olref[valpha]{defValphas} ينجح في تعريف
ما قصده. والمطلوب الأعم أن تصح التعريفات بالعودية المتناهية التجاوز؛
ولبيان ذلك عقدنا هذا القسم.

\emph{تنبيه: دقائق هذا البحث تحتاج إلى عناية}. أما أصله الجامع
فيسير: إن الاستقراء المتناهي التجاوز، منضمًا إلى الاستبدال، يصحح
صورًا عدة من العودية المتناهية التجاوز.
\footnote{يجوز للصيغ والحدود جميعًا أن تشتمل على معاملات، إلا حيث
ننص على خلاف ذلك.}

\begin{defn}
ليكن~$\tau({\OLId{latin-ordinary}{x}})$ حدًا، و${\OLId{latin-ordinary}{f}}$ دالة، و${\OLId{greek-variable}{alpha}}$ عددًا ترتيبيًا. نسمي
${\OLId{latin-ordinary}{f}}$ \emph{تقريب~${\OLId{greek-variable}{alpha}}$} لـ$\tau$ إذا وفقط إذا تحقق الشرطان
$\dom{{\OLId{latin-ordinary}{f}}} = {\OLId{greek-variable}{alpha}}$ و
$(\forall \beta \in {\OLId{greek-variable}{alpha}}){\OLId{latin-ordinary}{f}}(\beta) = \tau(\funrestrictionto{{\OLId{latin-ordinary}{f}}}{\beta})$.
\end{defn}

\begin{lem}[العودية المحدودة]\ollabel{transrecursionfun}
إذا أُعطي حد~$\tau({\OLId{latin-ordinary}{x}})$ وعدد ترتيبي~${\OLId{greek-variable}{alpha}}$، كان ثمة تقريب~${\OLId{greek-variable}{alpha}}$
واحد لا غير لـ$\tau$.
\end{lem}
\begin{proof}
نثبت، لكل~$\gamma \leq {\OLId{greek-variable}{alpha}}$، وجود تقريب~$\gamma$ ووحدانيته.

نبدأ بالوحدانية. خذ~${\OLId{latin-ordinary}{g}}$ و${\OLId{latin-ordinary}{h}}$ تقريبي~$\gamma$ و$\delta$، على الترتيب.
يقتضي الاستقراء المتناهي التجاوز على وسيطاتهما أن~${\OLId{latin-ordinary}{g}}(\beta) = {\OLId{latin-ordinary}{h}}(\beta)$
لكل~$\beta \in \dom{{\OLId{latin-ordinary}{g}}} \cap \dom{{\OLId{latin-ordinary}{h}}} = \gamma \cap \delta = \min(\gamma, \delta)$.
فالتقريبات، متى وجدت، وحيدة وتتفق حيث تشترك مجالاتها.

وأما الوجود فنثبته بالاستقراء البسيط المتناهي التجاوز
(\olref[ordinals][opps]{simpletransrecursion}) على الأعداد الترتيبية
$\delta \leq {\OLId{greek-variable}{alpha}}$.

الدالة الخالية تقريب~$\emptyset$ بلا ريب.

وإذا كانت~${\OLId{latin-ordinary}{g}}$ تقريب~$\gamma$، فإن~${\OLId{latin-ordinary}{g}} \cup \{\tuple{\gamma, \tau({\OLId{latin-ordinary}{g}})}\}$
تقريب~$\ordsucc{\gamma}$.

وأما إذا كان~$\gamma$ حديًا، وكانت~${\OLId{latin-ordinary}{g}}_\delta$ تقريب~$\delta$ لكل
$\delta < \gamma$، فنضع~${\OLId{latin-ordinary}{g}} = \bigcup_{\delta \in \gamma} {\OLId{latin-ordinary}{g}}_\delta$.
وهي دالة، لاتفاق دوال~${\OLId{latin-ordinary}{g}}_\delta$ في القيم حيث تتلاقى مجالاتها.
وإذا كان~$\delta \in \gamma$، وجدنا
${\OLId{latin-ordinary}{g}}(\delta) = {\OLId{latin-ordinary}{g}}_{\ordsucc{\delta}}(\delta) =
\tau(\funrestrictionto{{\OLId{latin-ordinary}{g}}_{\ordsucc{\delta}}}{\delta}) =
\tau(\funrestrictionto{{\OLId{latin-ordinary}{g}}}{\delta})$.

فتمت حالات الاستقراء المتناهي التجاوز جميعًا.
\end{proof}

وليس من اللازم أن نعرّف دالة؛ فإن اقتصرنا على تعريف \emph{حد} أمكن
رفع القيد~${\OLId{greek-variable}{alpha}}$ من النتيجة السابقة. وسنستعمل في القضية الآتية
وبرهانها، للحد~$\sigma$، الاصطلاح
$\funrestrictionto{\sigma}{{\OLId{greek-variable}{alpha}}} = \Setabs{\tuple{\beta, \sigma(\beta)}}{\beta \in {\OLId{greek-variable}{alpha}}}$.

\begin{thm}[العودية العامة]\ollabel{transrecursionschema}
لكل حد~$\tau({\OLId{latin-ordinary}{x}})$ يمكن تعريف حد~$\sigma({\OLId{latin-ordinary}{x}})$ تعريفًا صريحًا يحقق
$\sigma({\OLId{greek-variable}{alpha}}) = \tau(\funrestrictionto{\sigma}{{\OLId{greek-variable}{alpha}}})$ لكل عدد
ترتيبي~${\OLId{greek-variable}{alpha}}$.
\end{thm}

\begin{proof}
لكل~${\OLId{greek-variable}{alpha}}$، تضمن~\olref{transrecursionfun} تقريب~${\OLId{greek-variable}{alpha}}$ وحيدًا
${\OLId{latin-ordinary}{f}}_{\OLId{greek-variable}{alpha}}$ للحد~$\tau$. فلنعرّف~$\sigma({\OLId{greek-variable}{alpha}})$ بأنه
${\OLId{latin-ordinary}{f}}_{\ordsucc{{\OLId{greek-variable}{alpha}}}}({\OLId{greek-variable}{alpha}})$. وحينئذ:
\begin{align*}
\sigma({\OLId{greek-variable}{alpha}}) &=
{\OLId{latin-ordinary}{f}}_{\ordsucc{{\OLId{greek-variable}{alpha}}}}({\OLId{greek-variable}{alpha}}) \\&=
\tau(\funrestrictionto{{\OLId{latin-ordinary}{f}}_{\ordsucc{{\OLId{greek-variable}{alpha}}}}}{{\OLId{greek-variable}{alpha}}}) \\&=
\tau(\Setabs{\tuple{\beta, {\OLId{latin-ordinary}{f}}_{\ordsucc{{\OLId{greek-variable}{alpha}}}}(\beta)}}{\beta \in {\OLId{greek-variable}{alpha}}}) \\&=
\tau(\Setabs{\tuple{\beta, {\OLId{latin-ordinary}{f}}_{\ordsucc{\beta}}(\beta)}}{\beta \in {\OLId{greek-variable}{alpha}}})\\&=
\tau(\funrestrictionto{\sigma}{{\OLId{greek-variable}{alpha}}})
\end{align*}
وذلك لأن~${\OLId{latin-ordinary}{f}}_{\ordsucc{\beta}}(\beta) = {\OLId{latin-ordinary}{f}}_{\ordsucc{{\OLId{greek-variable}{alpha}}}}(\beta)$
لكل~$\beta < {\OLId{greek-variable}{alpha}}$، كما ثبت في~\olref{transrecursionfun}.
\end{proof}
\noindent
ولا بد من تذكر أن~\olref{transrecursionschema} \emph{مخطط}. فلا
نتوقع من~$\sigma$ أن تعرّف دالة، أي \emph{مجموعة} من نوع معين؛
وإلا كانت~$\dom{\sigma}$ مجموعة الأعداد الترتيبية كلها، وذلك محال
بمفارقة بورالي--فورتي (\olref[ordinals][basic]{buraliforti}).

وبقي وجه استناد تعريف مجموعات~${\OLId{latin-looped}{V}}_{\OLId{greek-variable}{alpha}}$ إلى~\olref{transrecursionschema}.
قد لا يظهر ذلك من أول الأمر؛ وتكشفه الصورة البسيطة الأخيرة من العودية
المتناهية التجاوز.

\begin{thm}[العودية البسيطة]\ollabel{simplerecursionschema}
متى أُعطي حدان~$\tau({\OLId{latin-ordinary}{x}})$ و$\theta({\OLId{latin-ordinary}{x}})$ ومجموعة~${\OLId{latin-looped}{A}}$، أمكن تعريف
حد~$\sigma({\OLId{latin-ordinary}{x}})$ صراحة بحيث:
\begin{align*}
\sigma(\emptyset) &= {\OLId{latin-looped}{A}}\\
\sigma(\ordsucc{{\OLId{greek-variable}{alpha}}}) &= \tau(\sigma({\OLId{greek-variable}{alpha}})) &&
\text{لكل عدد ترتيبي }{\OLId{greek-variable}{alpha}}\\
\sigma({\OLId{greek-variable}{alpha}}) &= \theta(\ran{\funrestrictionto{\sigma}{{\OLId{greek-variable}{alpha}}}})&&
\text{عندما يكون }{\OLId{greek-variable}{alpha}}\text{ عددًا ترتيبيًا حديًا}
\end{align*}
\end{thm}

\begin{proof}
نعرّف أولًا الحد~$\xi({\OLId{latin-ordinary}{x}})$ هكذا:
%t $\xi(x) = A$ if $x$ is not a function whose domain is an ordinal;
%otherwise let:
\[
\xi({\OLId{latin-ordinary}{x}}) =
\begin{cases}
{\OLId{latin-looped}{A}} &\text{إذا كان ${\OLId{latin-ordinary}{x}} = \emptyset$، أو لم يكن ${\OLId{latin-ordinary}{x}}$ دالة}\\
&\hspace{1em}\text{مجالها عدد ترتيبي؛ وخلاف ذلك:}\\
\tau({\OLId{latin-ordinary}{x}}({\OLId{greek-variable}{alpha}})) & \text{إذا كان $\dom{{\OLId{latin-ordinary}{x}}} = \ordsucc{{\OLId{greek-variable}{alpha}}}$}\\
\theta(\ran{{\OLId{latin-ordinary}{x}}}) & \text{إذا كان $\dom{{\OLId{latin-ordinary}{x}}}$ عددًا ترتيبيًا حديًا}
\end{cases}
\]
فبمقتضى~\olref{transrecursionschema} يوجد حد~$\sigma({\OLId{latin-ordinary}{x}})$ بحيث
$\sigma({\OLId{greek-variable}{alpha}}) = \xi(\funrestrictionto{\sigma}{{\OLId{greek-variable}{alpha}}})$ لكل عدد
ترتيبي~${\OLId{greek-variable}{alpha}}$. ثم إن~$\funrestrictionto{\sigma}{{\OLId{greek-variable}{alpha}}}$ دالة
مجالها~${\OLId{greek-variable}{alpha}}$. ونثبت الشروط المطلوبة لـ$\sigma$ بالاستقراء البسيط
المتناهي التجاوز (\olref[ordinals][opps]{simpletransrecursion}).

في حالة الأساس~$\sigma(\emptyset) = \xi(\emptyset) = {\OLId{latin-looped}{A}}$.

وفي حالة الخلف:
$\sigma(\ordsucc{{\OLId{greek-variable}{alpha}}}) =
\xi(\funrestrictionto{\sigma}{\ordsucc{{\OLId{greek-variable}{alpha}}}}) =
\tau(\funrestrictionto{\sigma}{\ordsucc{{\OLId{greek-variable}{alpha}}}}({\OLId{greek-variable}{alpha}})) =
\tau(\sigma({\OLId{greek-variable}{alpha}}))$.

وفي الحالة الأخيرة~$\sigma({\OLId{greek-variable}{alpha}}) = \xi(\funrestrictionto{\sigma}{{\OLId{greek-variable}{alpha}}}) =
\theta(\ran{\funrestrictionto{\sigma}{{\OLId{greek-variable}{alpha}}}})$، إذا كان~${\OLId{greek-variable}{alpha}}$ حديًا.
\end{proof}
\noindent
ولتسويغ~\olref[valpha]{defValphas} نأخذ~${\OLId{latin-looped}{A}} = \emptyset$
و$\tau({\OLId{latin-ordinary}{x}}) = \Pow{{\OLId{latin-ordinary}{x}}}$ و$\theta({\OLId{latin-ordinary}{x}}) = \bigcup {\OLId{latin-ordinary}{x}}$؛ وبذلك يستقيم أخيرًا
تعريف مجموعات~${\OLId{latin-looped}{V}}_{\OLId{greek-variable}{alpha}}$.{}

\end{document}
